\documentclass[11pt,a4paper]{article}

\input{../preambule}

\begin{document}

\title{Semaine 7}
\author{}
\date{16/11 - 22/11 }

\maketitle

\begin{exercise}
Montrer que l'ensemble des sous-groupes de $(\mathbb Q, +)$ est indénombrable.
\end{exercise}

\begin{exercise}
Soit $G$ un groupe. Un sous-groupe maximal de $G$ 
est un sous-groupe de $G$ différent de $G$ qui est élément maximal pour l'ordre de l'inclusion.
\begin{questions}
\item Quels sont les sous-groupes maximaux de $(\mathbb Z, +)$ ?
\item Même question pour $(\mathbb Q, +)$
\end{questions}
\end{exercise}

\begin{exercise}
Soit $G$ un groupe fini tel que $\forall x\in G, x^2=e$.
Montrer que $|G|$ est une puissance de $2$.
\end{exercise}

\begin{exercise}
Soit $G$ un groupe fini, $A,B$ deux parties de $G$
telles que $|A|+|B|>|G|$. Montrer que $AB=G$.
\end{exercise}

\begin{exercise}
Soit $G$ un groupe avec un nombre fini de sous-groupes.
Montrer que $G$ est fini.
\end{exercise}

\begin{exercise}
Soit $G$ un groupe fini tel que $Aut(G)=\set{id_G}$. Montrer que $\forall x\in G, x^2=e$.
\end{exercise}

\begin{exercise}
Trouver tous les $p$ premiers impairs tels que
$$\forall x,y\in\N,p\mid x^2+y^2\iff p\mid x \text{ et } p\mid y$$
\end{exercise}

\begin{exercise}
Soit $p=3k+2$ un nombre premier.\\
Soit $A\subset (\Z/p\Z)^{\times}$.\\
Pour $x\in\intset{1,p-1}$, on note $B(x)=A\cap\set{\overline{ix}, i\in\intset{k+1,2k+1}}$.
\begin{questions}
\item Calculer $\displaystyle\sum_{x=1}^{p-1} |B(x)|$
\item Soit $C\subset \Z\setminus\set{0}$ non vide et fini. 
Montrer que $C$ contient une partie $D$ telle que $|D|\geq\frac{|C|}{3}$ et
$\forall x,y,z\in D, x+y\neq z$
\end{questions}
\end{exercise}

\begin{exercise}
Soit $(G, \cdot)$ un groupe fini et $f:G\to \C^*$ un morphisme de groupes.
Calculer $\displaystyle\sum_{x\in G} f(x)$
\end{exercise}

\begin{exercise}
Soit $q,p\in\N$ tel que $q^p-1$ soit premier.
Montrer que $q=2$ et que $p$ est premier.
\end{exercise}

\begin{exercise}
Soit $k$ un diviseur de $2^p-1$ avec p premier. Montrer que $k\equiv 1\mod p$.
\end{exercise}

\begin{exercise}
Soit $k\in\N$ déterminer les couples solutions $(a,b)\in\N^2$ de l'équation
$a^2+b^2=2^k$ 
\end{exercise}





\end{document}